\subsection{Proof of Theorem~\ref{jep104local038}}

As explained in Section~3.2, any ring of coefficients considered above appears in an $\ell$-modular triple $(\jepPubBbb{K},\jepPubBbb{O},\jepPubBbb{L})$ where $\jepPubBbb{K}$ is a finite extension of $\jepPubBbb{Q}_\ell$, $\jepPubBbb{O}$ is its ring of integers, and $\jepPubBbb{L}$ is the residue field of $\jepPubBbb{O}$. Therefore we fix such a triple, and will treat the three cases in parallel.

The starting point of our proof will be the \emph{geometric Casselman--Shalika formula}, first conjectured in~\cite{jep104ref25} and then proved independently in~\cite{jep104ref26} and~\cite{jep104ref38} (see also Remark~\ref{jep104local025}\eqref{jep104local013}). We consider the composition
\[
 \chi_{U^+_{\jepPubScript{K}}} : U^+_{\jepPubScript{K}} \xrightarrow{\textstyle\chi_{\jepPubScript{K}}} (\jepPubBbb{G}_{\mathrm{a}})_{\jepPubScript{K}} \to \jepPubBbb{G}_{\mathrm{a}},
\]
where the second map is the ``residue'' morphism defined by
\[
 \sum_{i \in \jepPubBbb{Z}} f_i z^i \mapsto f_{-1}.
\]
For $\mu \in \boldsymbol{X}^\vee_+$ we set $S_\mu := U^+_{\jepPubScript{K}} \cdot L_\mu$; then there exists a unique function $\chi_\mu : S_\mu \to \jepPubBbb{G}_{\mathrm{a}}$ such that $\chi_\mu(u \cdot L_\mu)=\chi_{U^+_{\jepPubScript{K}}}(u)$ for any $u \in U^+_{\jepPubScript{K}}$. The geometric Casselman--Shalika formula states that for $\lambda,\mu \in \boldsymbol{X}^\vee_+$ we have
\begin{equation}
\label{jep104local002}
\mathsf{H}^i_c \bigl( S_\mu, \jepPubEuler{J}_{!*}(\lambda,\jepPubBbb{K})_{|S_\mu} \otimes_{\jepPubBbb{K}} \chi_\mu^*(\jepPubEuler{L}_\psi^{\jepPubBbb{K}}) \bigr) = \begin{cases}
\jepPubBbb{K} & \text{if $\lambda=\mu$ and $i=\langle 2\rho,\lambda\rangle$;} \\
0 & \text{otherwise.}
\end{cases}
\end{equation}

In the following lemma, we denote by $\chi_\mu' : z^{-\varsigma} X_{\mu+\varsigma} \to \jepPubBbb{G}_{\mathrm{a}}$ the unique function such that $\chi_\mu'(z^{-\varsigma} \cdot u \cdot L_{\mu+\varsigma}) = \chi_{I^+}(u)$ for $u \in I_\mathrm{u}^+$.

\begin{lem}
\label{jep104local022}
For $\Bbbk \in \{\jepPubBbb{K},\jepPubBbb{O},\jepPubBbb{L}\}$, for any $\lambda,\mu \in \boldsymbol{X}^\vee_+$ with $\lambda \neq \mu$ we have
\[
\mathsf{H}^{\langle \lambda+\mu,2\rho \rangle}_c \bigl( \mathrm{Gr}^\lambda \cap (z^{-\varsigma} X_{\mu+\varsigma}), (\chi_\mu')^*(\jepPubEuler{L}^\Bbbk_\psi)_{|\mathrm{Gr}^\lambda \cap (z^{-\varsigma} X_{\mu+\varsigma})}) = 0.
\]
\end{lem}

\begin{proof}
For $\alpha \in \Delta$ and $n \in \jepPubBbb{Z}_{\geqslant 0}$ we denote by $U_{\alpha,n} \subset G_{\jepPubScript{O}}$ the image of the morphism $x \mapsto u_\alpha(xz^n)$.
As explained e.g.~in~\cite[Lem.~2.2]{jep104ref38}, the action on $L_{\mu+\varsigma}$ induces an isomorphism
\[
\prod_{\alpha \in \Delta^+} \prod_{j=0}^{\langle \mu+\varsigma, \alpha \rangle -1} U_{\alpha,j} \xrightarrow{\sim} X_{\mu + \varsigma}.
\]
Multiplying by $z^{-\varsigma}$ we deduce that $z^{-\varsigma} X_{\mu+\varsigma} \subset S_\mu$, and moreover that $\chi_\mu'$ is the restriction of $\chi_\mu$ to $z^{-\varsigma} X_{\mu+\varsigma}$. By~\cite[Th.~3.2]{jep104ref36}, we have $\dim(\mathrm{Gr}^\lambda \cap S_\mu)=\langle \lambda+\mu, \rho \rangle$; it follows that $\dim(\mathrm{Gr}^\lambda \cap (z^{-\varsigma} X_{\mu+\varsigma})) \leqslant \langle \lambda+\mu, \rho \rangle$. If this inequality is strict, then our vanishing claim is obvious (see e.g.~\cite[Th.~I.8.8]{jep104ref23}). Otherwise, each irreducible component of $\mathrm{Gr}^\lambda \cap (z^{-\varsigma} X_{\mu+\varsigma})$ of dimension $\langle \lambda+\mu, \rho \rangle$ is dense (hence open) in an irreducible component of $\mathrm{Gr}^\lambda \cap S_\mu$; therefore to prove the lemma it suffices to prove that
\begin{equation}
\label{jep104local010}
\mathsf{H}^{\langle \lambda+\mu,2\rho \rangle}_c \bigl( \mathrm{Gr}^\lambda \cap S_\mu, \chi_\mu^*(\jepPubEuler{L}^\Bbbk_\psi)_{|\mathrm{Gr}^\lambda \cap S_\mu}) = 0.
\end{equation}
Finally, we note that since we are considering the top cohomology, the $\jepPubBbb{O}$-module $\mathsf{H}^{\langle \lambda+\mu,2\rho \rangle}_c \bigl( \mathrm{Gr}^\lambda \cap S_\mu, \chi_\mu^*(\jepPubEuler{L}^{\jepPubBbb{O}}_\psi)_{|\mathrm{Gr}^\lambda \cap S_\mu})$ is free, and the natural morphisms
\begin{align*}
\jepPubBbb{K} \otimes_{\jepPubBbb{O}} \mathsf{H}^{\langle \lambda+\mu,2\rho \rangle}_c \bigl( \mathrm{Gr}^\lambda \cap S_\mu, \chi_\mu^*(\jepPubEuler{L}^{\jepPubBbb{O}}_\psi)_{|\mathrm{Gr}^\lambda \cap S_\mu}) &\to \mathsf{H}^{\langle \lambda+\mu,2\rho \rangle}_c \bigl( \mathrm{Gr}^\lambda \cap S_\mu, \chi_\mu^*(\jepPubEuler{L}^{\jepPubBbb{K}}_\psi)_{|\mathrm{Gr}^\lambda \cap S_\mu}), \\
\jepPubBbb{L} \otimes_{\jepPubBbb{O}} \mathsf{H}^{\langle \lambda+\mu,2\rho \rangle}_c \bigl( \mathrm{Gr}^\lambda \cap S_\mu, \chi_\mu^*(\jepPubEuler{L}^{\jepPubBbb{O}}_\psi)_{|\mathrm{Gr}^\lambda \cap S_\mu}) &\to \mathsf{H}^{\langle \lambda+\mu,2\rho \rangle}_c \bigl( \mathrm{Gr}^\lambda \cap S_\mu, \chi_\mu^*(\jepPubEuler{L}^{\jepPubBbb{L}}_\psi)_{|\mathrm{Gr}^\lambda \cap S_\mu})
\end{align*}
are isomorphisms; hence it suffices to prove~\eqref{jep104local010} in case $\Bbbk=\jepPubBbb{K}$.

So, from now on we assume that $\Bbbk=\jepPubBbb{K}$. The geometric Casselman--Shalika formula~\eqref{jep104local002} implies that for any $\jepPubEuler{F}$ in $\mathrm{Perv}_{G_{\jepPubScript{O}}}(\mathrm{Gr},\jepPubBbb{K})$ we have $\mathsf{H}^i_c(S_\mu, \jepPubEuler{F}_{|S_\mu} \otimes_{\jepPubBbb{K}} \chi_\mu^*(\jepPubEuler{L}_\psi^{\jepPubBbb{K}}))=0$ for $i \neq \langle 2\rho,\mu \rangle$; therefore the morphism $(j_\lambda)_! \underline{\jepPubBbb{K}}_{\mathrm{Gr}^\lambda}[\langle \lambda, 2\rho \rangle] \to \jepPubEuler{J}_!(\lambda,\jepPubBbb{K})$ induces an isomorphism
\[
\mathsf{H}^{\langle \lambda+\mu,2\rho \rangle}_c \Bigl( S_\mu, \bigl((j_\lambda)_! \underline{\jepPubBbb{K}}_{\mathrm{Gr}^\lambda} \bigr)_{|S_\mu} \otimes_{\jepPubBbb{K}} \chi_\mu^*(\jepPubEuler{L}_\psi^{\jepPubBbb{K}}) \Bigr) \xrightarrow{\sim} \mathsf{H}^{\langle \mu,2\rho \rangle}_c \bigl( S_\mu, \jepPubEuler{J}_!(\lambda,\jepPubBbb{K}) \otimes_{\jepPubBbb{K}} \chi_\mu^*(\jepPubEuler{L}_\psi^{\jepPubBbb{K}}) \bigr).
\]
Now we have $\jepPubEuler{J}_!(\lambda,\jepPubBbb{K}) \cong \jepPubEuler{J}_{!*}(\lambda,\jepPubBbb{K})$ by Remark~\ref{jep104local024}\eqref{jep104local014}; hence the right-hand side vanishes if $\lambda \neq \mu$ by~\eqref{jep104local002}. On the other hand, the base change theorem shows that the left-hand side identifies with $\mathsf{H}^{\langle \lambda+\mu,2\rho \rangle}_c \bigl( \mathrm{Gr}^\lambda \cap S_\mu, \chi_\mu^*(\jepPubEuler{L}^{\jepPubBbb{K}}_\psi)_{|\mathrm{Gr}^\lambda \cap S_\mu})$; we have therefore proved~\eqref{jep104local010} in this case, hence the lemma.
\end{proof}

\begin{prop}
\label{jep104local023}
For $\Bbbk \in \{\jepPubBbb{K},\jepPubBbb{L}\}$, for any $\lambda,\mu \in \boldsymbol{X}^\vee_+$ with $\lambda \neq \mu$ we have
\[
\operatorname{Hom}_{\mathrm{Perv}_{\jepPubEuler{IW}}(\mathrm{Gr},\Bbbk)} \bigl( \Phi^0(\jepPubEuler{J}_!(\lambda,\Bbbk)), \nabla^{\jepPubEuler{IW}}_{\mu+\varsigma}(\Bbbk) \bigr)=0.
\]
\end{prop}

\begin{proof}
First, by exactness of $\Phi$ we see that the morphism $(j_\lambda)_!\underline{\Bbbk}_{\mathrm{Gr}^\lambda}[\langle\lambda,2\rho\rangle]\to\jepPubEuler{J}_!(\lambda,\Bbbk)$ induces an isomorphism
\[
\operatorname{Hom}_{D^{\mathrm{b}}_{\jepPubEuler{IW}}(\mathrm{Gr},\Bbbk)}\bigl(\Phi((j_\lambda)_!\underline{\Bbbk}_{\mathrm{Gr}^\lambda}[\langle\lambda,2\rho\rangle]),\nabla^{\jepPubEuler{IW}}_{\mu+\varsigma}(\Bbbk)\bigr)
\xrightarrow{\sim}\operatorname{Hom}_{\mathrm{Perv}_{\jepPubEuler{IW}}(\mathrm{Gr},\Bbbk)}\bigl(\Phi^0(\jepPubEuler{J}_!(\lambda,\Bbbk)),\nabla^{\jepPubEuler{IW}}_{\mu+\varsigma}(\Bbbk)\bigr).
\]
Let $J$ be the stabilizer of the point $L_\varsigma$ in $I^+_\mathrm{u}$. Then $\chi_{I^+}^*(\jepPubEuler{L}^\Bbbk_\psi)$ is trivial on $J$, so that we have a forgetful functor
\[
D^{\mathrm{b}}_{\jepPubEuler{IW}}(\mathrm{Gr},\Bbbk)\to D^{\mathrm{b}}_J(\mathrm{Gr},\Bbbk).
\]
By the same considerations as in Section~2.3, this functor admits a left adjoint, denoted ${}^{!}\mathsf{Ind}_J^{(I^+_\mathrm{u},\chi_{I^+})}$. Moreover, since $J$ is pro-unipotent the forgetful functor $D^{\mathrm{b}}_J(\mathrm{Gr},\Bbbk)\to D^{\mathrm{b}}_c(\mathrm{Gr},\Bbbk)$ is fully faithful.

If we denote by $\mathsf{F}_\varsigma$ the direct image under the automorphism $x\mapsto z^\varsigma\cdot x$ of $\mathrm{Gr}$, then from the definition we see that
\begin{equation}\label{jep104local005}
\Phi(\jepPubEuler{F})={}^{!}\mathsf{Ind}_J^{(I^+_\mathrm{u},\chi_{I^+})}\circ\mathsf{F}_\varsigma\circ\mathsf{For}_{z^{-\varsigma}Jz^\varsigma}^{G_{\jepPubScript{O}}}(\jepPubEuler{F})[-\langle\varsigma,2\rho\rangle]
\end{equation}
for any $\jepPubEuler{F}$ in $D^{\mathrm{b}}_{G_{\jepPubScript{O}}}(\mathrm{Gr},\Bbbk)$. In the setting of the proposition, we deduce an isomorphism
\[
\operatorname{Hom}_{D^{\mathrm{b}}_{\jepPubEuler{IW}}(\mathrm{Gr},\Bbbk)}\bigl(\Phi((j_\lambda)_!\underline{\Bbbk}_{\mathrm{Gr}^\lambda}[\langle\lambda,2\rho\rangle]),\nabla^{\jepPubEuler{IW}}_{\mu+\varsigma}(\Bbbk)\bigr)
\cong\operatorname{Hom}_{D^{\mathrm{b}}_c(\mathrm{Gr},\Bbbk)}\bigl((j_\lambda)_!\underline{\Bbbk}_{\mathrm{Gr}^\lambda}[\langle\lambda,2\rho\rangle],\mathsf{F}_\varsigma^{-1}(\nabla^{\jepPubEuler{IW}}_{\mu+\varsigma}(\Bbbk))[\langle\varsigma,2\rho\rangle]\bigr).
\]
Now $\mathsf{F}_\varsigma^{-1}(\nabla^{\jepPubEuler{IW}}_{\mu+\varsigma}(\Bbbk))$ identifies with the $*$-pushforward of $(\chi'_\mu)^*(\jepPubEuler{L}^\Bbbk_\psi)[\langle\mu+\varsigma,2\rho\rangle]$ under the embedding $z^{-\varsigma}X_{\mu+\varsigma}\to\mathrm{Gr}$. Hence, by the base change theorem, the right-hand side identifies with
\[
\mathsf{H}^{\langle\mu+2\varsigma-\lambda,2\rho\rangle}(\mathrm{Gr}^\lambda\cap z^{-\varsigma}X_{\mu+\varsigma},a^!(\chi'_\mu)^*(\jepPubEuler{L}^\Bbbk_\psi)),
\]
where $a:\mathrm{Gr}^\lambda\cap z^{-\varsigma}X_{\mu+\varsigma}\hookrightarrow z^{-\varsigma}X_{\mu+\varsigma}$ is the embedding.

So, we now need to show that $\mathsf{H}^{\langle\mu+2\varsigma-\lambda,2\rho\rangle}(\mathrm{Gr}^\lambda\cap z^{-\varsigma}X_{\mu+\varsigma},a^!(\chi'_\mu)^*(\jepPubEuler{L}^\Bbbk_\psi))$ vanishes. If $b$ denotes the unique map $z^{-\varsigma}X_{\mu+\varsigma}\to\mathrm{pt}$, then by Verdier duality we have
\begin{align*}
&\mathsf{H}^{\langle\mu+2\varsigma-\lambda,2\rho\rangle}(\mathrm{Gr}^\lambda\cap z^{-\varsigma}X_{\mu+\varsigma},a^!(\chi'_\mu)^*(\jepPubEuler{L}^\Bbbk_\psi))^*
=\mathsf{H}^{\langle\mu+2\varsigma-\lambda,2\rho\rangle}(b_*a^!(\chi'_\mu)^*(\jepPubEuler{L}^\Bbbk_\psi))^*\\
&\cong\mathsf{H}^{\langle\lambda-2\varsigma-\mu,2\rho\rangle}(b_!a^*\jepPubBbb{D}_{z^{-\varsigma}X_{\mu+\varsigma}}((\chi'_\mu)^*(\jepPubEuler{L}^\Bbbk_\psi))).
\end{align*}
Now since $z^{-\varsigma}X_{\mu+\varsigma}$ is smooth of dimension $\langle\mu+\varsigma,2\rho\rangle$ we have an isomorphism
$\jepPubBbb{D}_{z^{-\varsigma}X_{\mu+\varsigma}}((\chi'_\mu)^*(\jepPubEuler{L}^\Bbbk_\psi))\cong(\chi'_\mu)^*(\jepPubEuler{L}^\Bbbk_{-\psi})[2\langle\mu+\varsigma,2\rho\rangle]$, which shows that
\[
\mathsf{H}^{\langle\mu+2\varsigma-\lambda,2\rho\rangle}(\mathrm{Gr}^\lambda\cap z^{-\varsigma}X_{\mu+\varsigma},a^!(\chi'_\mu)^*(\jepPubEuler{L}^\Bbbk_\psi))^*
\cong\mathsf{H}^{\langle\mu+\lambda,2\rho\rangle}_c(\mathrm{Gr}^\lambda\cap z^{-\varsigma}X_{\mu+\varsigma},a^*(\chi'_\mu)^*(\jepPubEuler{L}^\Bbbk_{-\psi})).
\]
The right-hand side vanishes by Lemma~\ref{jep104local022}, hence so does the left-hand side, which completes the proof.
\end{proof}

We can finally give the proof of Theorem~\ref{jep104local038}.

\begin{proof}[Proof of Theorem~{\rm \ref{jep104local038}}]
Let $\tau : G_{\jepPubScript{K}} \to \mathrm{Gr}$ be the projection. Let $\lambda \in \boldsymbol{X}^\vee_+$, and denote by $m_{\varsigma,\lambda}$ the restriction of $m^{\mathrm{Gr}}$ to $\tau^{-1}(\overline{\strut \mathrm{Gr}^\varsigma}) \times^{G_{\jepPubScript{O}}} \overline{\mathrm{Gr}^\lambda}$ (see Section~2.4 for the notation). Then it is well known that:
\begin{itemize}
\item
$m_{\varsigma,\lambda}$ takes values in $\overline{\mathrm{Gr}^{\lambda+\varsigma}} = \overline{\strut X_{\lambda+\varsigma}}$;
\item
its restriction to the preimage of $X_{\lambda+\varsigma}$ is an isomorphism;
\item
this preimage is contained in $\tau^{-1}(X_\varsigma) \times^{G_{\jepPubScript{O}}} \mathrm{Gr}^\lambda$.
\end{itemize}

These properties imply
that the perverse sheaf $\Phi^0(\jepPubEuler{J}_!(\lambda,\Bbbk))$ is supported on $\overline{X_{\lambda+\varsigma}}$, and that its restriction to $X_{\lambda+\varsigma}$ is a perversely shifted local system of rank $1$. The same comments apply to $\Phi^0(\jepPubEuler{J}_*(\lambda,\Bbbk))$.
Hence there exist canonical morphisms
\[
f^\Bbbk_\lambda : \Delta^{\jepPubEuler{IW}}_{\lambda + \varsigma}(\Bbbk) \to \Phi^0(\jepPubEuler{J}_!(\lambda,\Bbbk)) \quad \text{and} \quad g_\lambda^\Bbbk : \Phi^0(\jepPubEuler{J}_*(\lambda,\Bbbk)) \to \nabla^{\jepPubEuler{IW}}_{\lambda + \varsigma}(\Bbbk)
\]
whose restrictions to $X_{\lambda+\varsigma}$ are isomorphisms.

We claim that $f_\lambda^\Bbbk$ is an isomorphism.
First we note that all of our constructions are compatible with extension-of-scalars functors in the obvious sense (see in particular~\cite[Prop.~8.1]{jep104ref36} for the case of $\jepPubEuler{J}_!(\lambda,\Bbbk)$; the case of the Whittaker standard object is much easier since no perverse truncation is involved). If $\Bbbk \in \{\jepPubBbb{K},\jepPubBbb{L}\}$, by Proposition~\ref{jep104local023} we know that $\Phi^0(\jepPubEuler{J}_!(\lambda,\Bbbk))$ has no quotient of the form $\mathrm{IC}^{\jepPubEuler{IW}}_{\mu+\varsigma}(\Bbbk)$ with $\mu \neq \lambda$; therefore $f_\lambda^\Bbbk$ is surjective. 

The surjectivity of $f_\lambda^{\jepPubBbb{L}}$ implies that $f_\lambda^{\jepPubBbb{O}}$ must be surjective also. On the other hand, by Remark~\ref{jep104local025}\eqref{jep104local013} the object $\Delta^{\jepPubEuler{IW}}_{\lambda + \varsigma}(\jepPubBbb{K})$ is simple; hence $f_\lambda^{\jepPubBbb{K}}$ is injective, which implies that $\ker(f_\lambda^{\jepPubBbb{O}})$ is a torsion object. Since this object embeds in the torsion-free object $\Delta^{\jepPubEuler{IW}}_{\lambda + \varsigma}(\jepPubBbb{O})$, it must be zero. We finally obtain that $f_\lambda^{\jepPubBbb{O}}$ is an isomorphism, so that $f_\lambda^{\jepPubBbb{K}}$ and $f_\lambda^{\jepPubBbb{L}}$ are isomorphisms as well.

Once we know that $f_\lambda^\Bbbk$ is an isomorphism, by Verdier duality (see the comments preceding Lemma~\ref{jep104local016}) we deduce that $g_\lambda^\Bbbk$ is an isomorphism as well. (More precisely, we use the claim about $f_\lambda^\Bbbk$ in the setting where $\psi$ is replaced by $\psi^{-1}$, and the fact that $\jepPubBbb{D}_{\mathrm{Gr}}(\jepPubEuler{J}_!(\lambda,\Bbbk)) = \jepPubEuler{J}_*(\lambda,\Bbbk)$, see~\cite[Prop.~8.1(c)]{jep104ref36}.)

Now we
conclude the proof as follows.
Since $\Phi^0$ is exact, it induces a functor
\[
D^{\mathrm{b}}(\Phi^0) : D^{\mathrm{b}} \mathrm{Perv}_{G_{\jepPubScript{O}}}(\mathrm{Gr}, \Bbbk) \to D^{\mathrm{b}} \mathrm{Perv}_{\jepPubEuler{IW}}(\mathrm{Gr}, \Bbbk).
\]
We will prove that $D^{\mathrm{b}}(\Phi^0)$ is an equivalence, which will imply that $\Phi^0$ is an equivalence as well, hence will conclude the proof.
It is not difficult to see that the category $D^{\mathrm{b}} \mathrm{Perv}_{G_{\jepPubScript{O}}}(\mathrm{Gr}, \Bbbk)$, resp.~$D^{\mathrm{b}} \mathrm{Perv}_{\jepPubEuler{IW}}(\mathrm{Gr}, \Bbbk)$, is generated as a triangulated category by the objects $\{\jepPubEuler{J}_!(\lambda,\Bbbk) : \lambda \in \boldsymbol{X}^\vee_+\}$, resp.~by the objects $\{\Delta^{\jepPubEuler{IW}}_{\lambda+\varsigma}(\Bbbk) : \lambda \in \boldsymbol{X}^\vee_+\}$, as well as by the objects $\{\jepPubEuler{J}_*(\lambda,\Bbbk) : \lambda \in \boldsymbol{X}^\vee_+\}$, resp.~by the objects $\{\nabla^{\jepPubEuler{IW}}_{\lambda+\varsigma}(\Bbbk) : \lambda \in \boldsymbol{X}^\vee_+\}$. Hence to conclude it suffices to prove that for any $\lambda,\mu \in \boldsymbol{X}^\vee_+$ and any $n \in \jepPubBbb{Z}$ the functor $\Phi^0$ induces an isomorphism
\[
 \operatorname{Ext}^n_{\mathrm{Perv}_{G_{\jepPubScript{O}}}(\mathrm{Gr}, \Bbbk)}(\jepPubEuler{J}_!(\lambda,\Bbbk), \jepPubEuler{J}_*(\mu,\Bbbk)) \xrightarrow{\sim} \operatorname{Ext}^n_{\mathrm{Perv}_{\jepPubEuler{IW}}(\mathrm{Gr}, \Bbbk)}(\Delta^{\jepPubEuler{IW}}_{\lambda+\varsigma}(\Bbbk), \nabla^{\jepPubEuler{IW}}_{\mu+\varsigma}(\Bbbk)).
\]
However, this is clear from Remark~\ref{jep104local024}\eqref{jep104local011} and Remark~\ref{jep104local025}\eqref{jep104local012}.
\end{proof}

\begin{rmk}
\begin{enumerate}
\item
One can explicitly describe the inverse to $\Phi^0$, as follows. In view of~\eqref{jep104local005}, the functor
\[
\Psi :={}^{*}\mathsf{Ind}_{z^{-\varsigma} J z^{\varsigma}}^{G_{\jepPubScript{O}}} \circ \mathsf{F}_\varsigma^{-1} \circ \mathsf{For}^{(I_\mathrm{u}^+, \chi_{I^+})}_J [\langle \varsigma,2\rho \rangle] : D^{\mathrm{b}}_{\jepPubEuler{IW}}(\mathrm{Gr},\Bbbk) \to D^{\mathrm{b}}_{G_{\jepPubScript{O}}}(\mathrm{Gr},\Bbbk)
\]
is right adjoint to $\Phi$. Since $\Phi$ is exact, $\Psi$ is left exact, and the functor $\Psi^0:={}^{p}\jepPubEuler{H}^0 \circ \Psi_{|\mathrm{Perv}_{\jepPubEuler{IW}}(\mathrm{Gr},\Bbbk)}$ is right adjoint to $\Phi^0$. Since $\Phi^0$ is an equivalence, $\Psi^0$ must be its inverse.
\item
 From the point of view suggested by the Finkelberg--Mirkovi\'c conjecture (see~Section~\ref{jep104local036}), the isomorphisms $f^\Bbbk_\lambda$ and $g^\Bbbk_\lambda$ are geometric analogues of the isomorphism stated in~\cite[Prop.~II.3.19]{jep104ref29}.

\end{enumerate}
\end{rmk}

